\newpage
\chapter{Описание лабораторной установки}

Виртуальная лабораторная установка представляет собой интерактивный измерительный стенд, развернутый в среде MATLAB\@.
Стенд содержит верхний блок органов управления (ручки настроек) и пять измерительных панелей (экранов), имитирующих работу реальных лабораторных приборов.

Лабораторная установка включает в себя:
\begin{enumerate}
    \item Блок ввода параметров (генератор): выбор типа модуляции (из 12 возможных), амплитуда $U_0$ (В), несущая частота $f_0$ (кГц), индекс модуляции $m$.
    \item Панель В: Экран НЧ первичного сообщения в виде синусоиды или бинарного сигнала.
    \item Панель А: Цифровой осциллограф результирующего ВЧ-радиосигнала во времени.
    \item Панель Б: Анализатор спектра (БПФ) --- частотная область в линейном масштабе (В) или в логарифмическом масштабе (дБмкВ).
    \item Панель Г: Векторный анализатор --- фазовый портрет сигнала на комплексной IQ-плоскости.
    \item Панель Д: Измерительное табло параметров: $U_m$, $f_0$, $m$, девиация частоты или фазы, скорость.
\end{enumerate}